\documentclass[aps,preprint]{revtex4}%
\usepackage{amsfonts}
\usepackage{amsmath}
\usepackage{amssymb}
\usepackage{graphicx}%
\setcounter{MaxMatrixCols}{30}
%TCIDATA{OutputFilter=latex2.dll}
%TCIDATA{Version=4.10.0.2347}
%TCIDATA{CSTFile=revtex4.cst}
%TCIDATA{Created=Sunday, December 01, 2002 18:14:35}
%TCIDATA{LastRevised=Tuesday, December 17, 2002 12:54:24}
%TCIDATA{<META NAME="GraphicsSave" CONTENT="32">}
%TCIDATA{<META NAME="DocumentShell" CONTENT="Articles\SW\REVTeX 4">}
\newtheorem{theorem}{Theorem}
\newtheorem{acknowledgement}[theorem]{Acknowledgement}
\newtheorem{algorithm}[theorem]{Algorithm}
\newtheorem{axiom}[theorem]{Axiom}
\newtheorem{claim}[theorem]{Claim}
\newtheorem{conclusion}[theorem]{Conclusion}
\newtheorem{condition}[theorem]{Condition}
\newtheorem{conjecture}[theorem]{Conjecture}
\newtheorem{corollary}[theorem]{Corollary}
\newtheorem{criterion}[theorem]{Criterion}
\newtheorem{definition}[theorem]{Definition}
\newtheorem{example}[theorem]{Example}
\newtheorem{exercise}[theorem]{Exercise}
\newtheorem{lemma}[theorem]{Lemma}
\newtheorem{notation}[theorem]{Notation}
\newtheorem{problem}[theorem]{Problem}
\newtheorem{proposition}[theorem]{Proposition}
\newtheorem{remark}[theorem]{Remark}
\newtheorem{solution}[theorem]{Solution}
\newtheorem{summary}[theorem]{Summary}
\newenvironment{proof}[1][Proof]{\noindent\textbf{#1.} }{\ \rule{0.5em}{0.5em}}
\begin{document}
\preprint{HEP/123-qed}
\title[Short title for running header]{Long title that appears on the title page}
\author{First Author}
\affiliation{My Institution}
\author{Second Author}
\affiliation{My Institution}
\author{Third Author}
\affiliation{Other Institution}
\keywords{one two three}
\pacs{PACS number}

\begin{abstract}
Shell document for REV\TeX{} 4.

\end{abstract}
\volumeyear{year}
\volumenumber{number}
\issuenumber{number}
\eid{identifier}
\date[Date text]{date}
\received[Received text]{date}

\revised[Revised text]{date}

\accepted[Accepted text]{date}

\published[Published text]{date}

\startpage{101}
\endpage{102}
\maketitle
\tableofcontents


A $p$-particle operator component, $T_{p}^{N}$, of $T$ is in general a sum of
all possible forms of $p$-particle operators i.e. products and sums of lower
order ones as well as irreducible components. The concept of an irreducible
part of an operator is defined as one that cannot be expressed in terms of
lower order parts in the following way:

\begin{enumerate}
\item[1] We define a one particle part to be irreducible

\item[2] A two particle part is irreducible if it cannot be expressed in terms
of a one particle part as
\begin{equation}
T_{2}\left(  \mathsf{z}_{1},\mathsf{z}_{2}\right)  \neq T_{1}\left(
\mathsf{z}_{1}\right)  T_{1}\left(  \mathsf{z}_{2}\right)
\end{equation}
nor
\begin{equation}
T_{2}\left(  \mathsf{z}_{1},\mathsf{z}_{2}\right)  \neq T_{1}\left(
\mathsf{z}_{1}\right)  \delta\left(  \mathsf{z}_{2}\right)  +\delta\left(
\mathsf{z}_{1}\right)  T_{1}\left(  \mathsf{z}_{2}\right)
\end{equation}

\end{enumerate}

In general a $p$-particle part is irreducible if it cannot be expressed in
terms of a $p-1$ parts as
\begin{equation}
T_{p}\left(  \mathsf{z}_{1},\cdots,\mathsf{z}_{p}\right)  \neq T_{q}\left(
\mathsf{z}_{1},\cdots,\mathsf{z}_{q}\right)  \cdots T_{q}\left(
\mathsf{z}_{p-q},\cdots,\mathsf{z}_{p}\right)  \text{ for some }q
\end{equation}
nor
\begin{equation}
T_{p}\left(  \mathsf{z}_{\sigma\left(  1\right)  },\cdots,\mathsf{z}%
_{\sigma\left(  p\right)  }\right)  \neq%
%TCIMACRO{\dsum \limits_{\sigma\in\mathcal{S}_{p}}}%
%BeginExpansion
{\displaystyle\sum\limits_{\sigma\in\mathcal{S}_{p}}}
%EndExpansion
T_{p-1}\left(  \mathsf{z}_{\sigma\left(  1\right)  }\cdots\mathsf{z}%
_{\sigma\left(  p-1\right)  }\right)  \delta\left(  \mathsf{z}_{\sigma\left(
p\right)  }\right)
\end{equation}
A component is irreducible if it is composed of irreducible parts e.g.
$T_{p}^{N}$ is irreducible where
\begin{equation}
T_{p}^{N}=%
%TCIMACRO{\dsum \limits_{\sigma\in\mathcal{S}_{N}}}%
%BeginExpansion
{\displaystyle\sum\limits_{\sigma\in\mathcal{S}_{N}}}
%EndExpansion
T_{p}\left(  \mathsf{z}_{\sigma\left(  1\right)  }\cdots\mathsf{z}%
_{\sigma\left(  p\right)  }\right)  \delta\left(  \mathsf{z}_{\sigma\left(
N-p\right)  }\right)
\end{equation}
if
\begin{equation}
T_{p}\left(  \mathsf{z}_{1},\cdots,\mathsf{z}_{p}\right)  \neq T_{q}\left(
\mathsf{z}_{1},\cdots,\mathsf{z}_{q}\right)  \cdots T_{q}\left(
\mathsf{z}_{p-q},\cdots,\mathsf{z}_{p}\right)  \text{ for some }q
\end{equation}
nor%
\begin{equation}
T_{p}\left(  \mathsf{z}_{\sigma\left(  1\right)  },\cdots,\mathsf{z}%
_{\sigma\left(  p\right)  }\right)  \neq%
%TCIMACRO{\dsum \limits_{\sigma\in\mathcal{S}_{p}}}%
%BeginExpansion
{\displaystyle\sum\limits_{\sigma\in\mathcal{S}_{p}}}
%EndExpansion
T_{p-1}\left(  \mathsf{z}_{\sigma\left(  1\right)  }\cdots\mathsf{z}%
_{\sigma\left(  p-1\right)  }\right)  \delta\left(  \mathsf{z}_{\sigma\left(
p\right)  }\right)  .
\end{equation}


From these classes we can also define subclasses of the form $\left\{
T_{n_{1}}\left(  \mathsf{z}_{1}^{n_{1}}\right)  \vee\delta^{N-n_{1}}\left(
\mathsf{z}_{n_{1}+1}^{N-n_{1}}\right)  \right\}  $ for $1\leq n_{1}\leq N-1$
where $T_{n_{j}}\left(  \mathsf{z}_{j}^{n_{j}}\right)  \in F^{n_{j}}$.

Developing eq. $\left(  \ref{gso}\right)  $ we get
\begin{align}
&
%TCIMACRO{\dint }%
%BeginExpansion
{\displaystyle\int}
%EndExpansion
\left(  \Phi\left(  \mathsf{z}_{1},\mathsf{z}^{N-1}\right)  \Phi\left(
\mathsf{z}_{1}^{\prime},\mathsf{z}^{N-1}\right)  ^{\ast}+\Phi\left(
\mathsf{z}_{1},\mathsf{z}^{N-1}\right)  +\Phi\left(  \mathsf{z}_{1}^{\prime
},\mathsf{z}^{N-1}\right)  ^{\ast}\right)  D_{0_{F}}^{N}\left[  \rho\right]
\left(  \mathsf{z}_{1},\mathsf{z}^{N-1},\mathsf{z}_{1}^{\prime},\mathsf{z}%
^{N-1}\right)  d\mathsf{z}^{N-1}\overset{a.e}{=}0\nonumber\\
&  \Rightarrow\int\Phi\left(  \mathsf{z}_{1},\mathsf{z}^{N-1}\right)
D_{0_{F}}^{N}\left[  \rho\right]  \left(  \mathsf{z}_{1},\mathsf{z}%
^{N-1},\mathsf{z}_{1}^{\prime},\mathsf{z}^{N-1}\right)  \Phi\left(
\mathsf{z}_{1}^{\prime},\mathsf{z}^{N-1}\right)  ^{\ast}d\mathsf{z}%
^{N-1}\nonumber\\
&  \qquad\qquad\qquad\qquad\qquad\left.  =\right.  -\int\left\{  \Phi\left(
\mathsf{z}_{1},\mathsf{z}^{N-1}\right)  D_{0_{F}}^{N}\left[  \rho\right]
\left(  \mathsf{z}_{1},\mathsf{z}^{N-1},\mathsf{z}_{1}^{\prime},\mathsf{z}%
^{N-1}\right)  +D_{0_{F}}^{N}\left[  \rho\right]  \left(  \mathsf{z}%
_{1},\mathsf{z}^{N-1},\mathsf{z}_{1}^{\prime},\mathsf{z}^{N-1}\right)
\Phi\left(  \mathsf{z}_{1}^{\prime},\mathsf{z}^{N-1}\right)  ^{\ast}\right\}
d\mathsf{z}^{N-1}%
\end{align}


\bigskip%

\begin{align*}
&  -\left\langle \left(  \mathbf{\nabla}_{j}-\mathbf{\nabla}_{j}\omega\right)
\nu_{N}\Psi|\left(  \mathbf{\nabla}_{j}-\mathbf{\nabla}_{j}\omega\right)
\nu_{N}\Psi\right\rangle \\
&  =\left\langle \left\{  \left(  \mathbf{\nabla}_{j}\nu_{N}\right)  +\nu
_{N}\mathbf{\nabla}_{j}-\nu_{N}\left(  \mathbf{\nabla}_{j}\omega_{N}\right)
\right\}  \Psi|\left\{  \left(  \mathbf{\nabla}_{j}\nu_{N}\right)  +\nu
_{N}\mathbf{\nabla}_{j}-\nu_{N}\left(  \mathbf{\nabla}_{j}\omega_{N}\right)
\right\}  \Psi\right\rangle \\
&  =\left\langle \Psi|\left\{  \mathbf{\nabla}_{j}\nu_{N}\bullet
\mathbf{\nabla}_{j}\nu_{N}+\nu_{N}\left(  \mathbf{\nabla}_{j}\nu_{N}\right)
\bullet\mathbf{\nabla}_{j}-\nu_{N}\mathbf{\nabla}_{j}\nu_{N}\bullet
\mathbf{\nabla}_{j}\omega\right.  \right.  \\
&  +\nu_{N}\left(  \mathbf{\nabla}_{j}\nu_{N}\right)  \bullet\mathbf{\nabla
}_{j}+\left(  \nu_{N}\right)  ^{2}\mathbf{\nabla}_{j}^{2}-\nu_{N}%
\mathbf{\nabla}_{j}\nu_{N}\bullet\mathbf{\nabla}_{j}\omega-\left(  \nu
_{N}\right)  ^{2}\left(  \mathbf{\nabla}_{j}^{2}\omega\right)  -\nu
_{N}\mathbf{\nabla}_{j}\omega\bullet\mathbf{\nabla}_{j}\\
&  \left.  -\nu_{N}\mathbf{\nabla}_{j}\omega_{N}\bullet\mathbf{\nabla}_{j}%
\nu_{N}-\left(  \nu_{N}\right)  ^{2}\mathbf{\nabla}_{j}\omega_{N}%
\bullet\mathbf{\nabla}_{j}+\nu_{N}^{2}\mathbf{\nabla}_{j}\omega_{N}%
\bullet\mathbf{\nabla}_{j}\omega_{N}\right\}  \left.  \Psi\right\rangle
\end{align*}



\end{document}